\section{Position Is the Whole Contract}
\label{sec:ch11-position-is-the-whole-contract}

\index{entities@\code{\_entities}!the order clause, twice|(}%
Chapter~\ref{ch:router-n-plus-1} quoted Apollo's subgraph specification saying
that the entity route
\enquote{must return a list of entity objects that correspond to the provided
representations, in the exact same order}, and that
\enquote{entries in the list can be null if no entity exists for a provided
representation}~\autocite{apollosubgraphspec}. That was the sentence pair I
needed there, and I read past a second statement of the same rule that matters
more here. Further down, the page gives a subgraph library an algorithm for
answering the field, and the step that puts an entity into the answer reads
\enquote{add the fetched entity object to the array of entity objects}, followed
immediately by \enquote{make sure objects are listed in the same order as their
corresponding representations}~\autocite{apollosubgraphspec}.

\index{entities@\code{\_entities}!stated twice, checked nowhere}%
Once as a contract on the answer, once as an instruction to the implementer.
I read the repetition as a warning, and I would rather take it as one than
find out why it is there.

\index{entities@\code{\_entities}!nothing validates the answer}%
What the page does not contain is any requirement to check. Its validation
rules are all about the representation coming in: that it carries
\code{\_\_typename}, and that it carries every field named in the
\code{@key} field set. About the answer it says two things, order and
nullability, and there is no third sentence asking anybody to confirm that the
entity at position two is the entity representation two named. I went looking
for one, in the same spirit as chapter~\ref{ch:router-n-plus-1}'s search for a
sentence about batching, and there is nothing.

\subsection{What the Router Does With the List}

\index{router!merges by position}%
The router's half of it is one loop, and it is worth being exact about,
because the shape of the failure follows from it. The executor inside the
router binary walks the answer array and the list of parent objects together,
merging each entry of one into the entry of the other that sits at the same
index. Where
chapter~\ref{ch:router-n-plus-1}'s de-duplication happened, it walks the map it
built from representation hash to batch position instead, and merges the answer
at that position into every parent that hashed to it.

\index{router!checks the length and nothing else}%
Before either loop it makes one check, and only one: that the number of entries
it got back equals the number of representations it sent. A short list or a
long one is refused, and the router renders an error saying the returned
entities count does not match the count of representation variables. That
check is why a subgraph cannot quietly drop an entry.

\index{router!a reordered list is not detectable}%
A list of the right length in the wrong order passes it without comment. There
is nothing else to catch it. The merge is a JSON deep merge with no notion of
identity: it does not read \code{\_\_typename} off the returned object, it does
not compare a key field, and it could not, because it did not ask for the key.
Look again at the entity fetch the planner wrote for the failing request:

\begin{minted}{graphql}
query($representations: [_Any!]!) {
  _entities(representations: $representations) {
    ... on Speaker { __typename name }
  }
}
\end{minted}

\index{router!never asks for the key it sent}%
The client asked for \code{name}, so the router asks for \code{name}. The
\code{id} is in the representation going out and is nowhere in the selection
set coming back. The router has no copy of the answer's identity to compare
against, and the reason is not an oversight: asking for a key the router
already knows would be a column nobody wanted, on every entity fetch in every
graph, forever.

Figure~\ref{fig:ch11-position-is-the-join} is the whole path.

\begin{figure}[htbp]
  \centering
  % The only thing joining an _entities answer to the question that asked for
% it is its position in the list, and three separate parties agree to that
% and none of them checks it: the router that built the representation list,
% the loader that turned rows back into a list, and the router again when it
% merges that list into the parent response.
%
% This is a positional-correspondence figure, not a timeline, and that is a
% deliberate departure from every figure since ch01-one-field.tex: the axis,
% tick, event, span and lane styles those files share all draw order as
% something that unfolds left to right over one request. Here the subject is
% not sequence, it is a list index read down four stages, so the picture is
% four rows of two boxes each, one box per list position, and a reader
% follows a single column straight down to see where the identity in it
% stops matching the label on it. The nearest relative in this book is
% ch04-null-climbs.tex, which is also a shape rather than a timeline; the
% box style, the link style and the drawn-sample legend are copied from
% there. What is NOT copied is what solid and dashed mean; see below.
%
% The four rows, top to bottom, and the measured data inside them (SPEC:
% do not alter these facts):
%   1. what the router asked: the client requested session 4 (The Bank That
%      Split Its Graph) then session 1 (Schemas That Outlive Their Authors),
%      so the representation list the router builds is position 0 =
%      Speaker:3 (Chidi Okafor, the speaker of session 4), position 1 =
%      Speaker:1 (Ada Fischer, the speaker of session 1).
%   2. what the store returned: one statement, which as StatementLog prints
%      it reads WHERE "s"."Id" IN (@ids1, @ids2) and carries keys 3 and 1
%      in that order, and the store answers in its own row order, ascending
%      id: row 0 is Ada Fischer (id 1), row 1 is Chidi Okafor (id 3). The
%      picture names the keys rather than writing SQL, because SQL with
%      literals in it is not what any capture in this book shows. This is
%      where the
%      order the router asked in and the order the store answered in stop
%      matching, and it is the only place in the whole figure that matters.
%   3. what the subgraph answered: the loader writes row 0 into slot 0 and
%      row 1 into slot 1 with no check against which key asked for which
%      row, so position 0 is now labeled Ada Fischer and position 1 Chidi
%      Okafor.
%   4. what the client was told: the router merges answer position 0 into
%      parent position 0 and answer position 1 into parent position 1, so
%      The Bank That Split Its Graph comes back credited to Ada Fischer and
%      Schemas That Outlive Their Authors to Chidi Okafor. Both are real
%      people in the graph, and the response is HTTP 200 with no error
%      anywhere in it.
%
% Solid, arrowed lines join row 1 to row 2 and are the only lines in the
% figure that cross: they show where the data for each requested identity
% actually landed, and the store handed it back in the opposite order from
% the one the router asked in. Dashed lines join row 2 to row 3 and row 3
% to row 4, drawn straight down with no crossing at all, because every step
% after the store answers is a position carried forward unchecked; nobody
% asks whether the entity in the slot is the one that was asked for.
%
% Solid versus dashed is NOT inherited from ch04-null-climbs.tex, and an
% earlier draft of this header wrongly said it was. There, dashed is a
% property of a BOX: a position that does not exist in the response at all.
% Here it is a property of a LINK: a correspondence nobody checked. The two
% figures are not in conflict, because neither uses the style for the other
% thing, but they are not the same convention and a later editor should not
% reconcile them. The legend at the foot of this picture exists precisely
% because the meaning is local to it. The arrowhead on the solid lines
% only is deliberate: it puts the one directional claim in the picture on
% the one relationship the figure is arguing for, and leaves the passive
% carry-forward lines without one.
%
% The row label at the top doubles as the column header: position 0 and
% position 1 are named once, in row 1, next to the representations that
% first fill them, rather than repeated in the boxes of every row, so a
% reader is not re-told the column number four times to learn it once.
%
% Span and spanlabel below the fourth row are the brace idiom from
% ch01-one-field.tex, ch09-two-fetches.tex and ch10-one-fetch-two-costs.tex,
% reused unchanged, carrying the one sentence the whole figure is building
% toward: nothing between the store and the client ever compares a returned
% entity with the representation that asked for it. The two-line legend
% below that is an idiom ch04-null-climbs.tex already uses, a drawn sample of each
% link style next to a one-line gloss, because this figure is the first one
% since ch04 to give solid and dashed a meaning that is not self-evident
% from an arrowhead and a timeline.
%
% Widths: two columns of boxes, centered at x=45 and x=105, each 40mm of text
% plus 2mm inner sep on either side, so the widest boxes span x=23 to x=67
% and x=83 to x=127. The bottom brace and its label stay inside x=24 to
% x=126 and x=18 to x=143 respectively. Row labels were once the widest text
% here and are not any more: they were cut to short tags because at their
% original length they ran straight through the crossing lines, and the one
% detail that had to survive, which keys the statement asked for, moved to
% its own right-anchored note at x=150. The longest row label is now 43
% characters, and ch10's own header measures a comparable line at about
% 1.4mm per character, so it ends near x=60, clear of the left column at
% x=45 only because the columns start below it rather than beside it.
% Everything in this figure therefore stays inside the 155mm measure this
% book's geometry allows (A4, 30mm inner, 25mm outer). Re-check the row
% labels against the crossing if any of them is ever lengthened again.
%
% No timings anywhere, per SPEC decision 19: the vertical position of a row
% is sequence only, not duration.
%
% No quotation marks anywhere in this file. The book's strict Ascii mode
% (see check-chapter.psd1) makes the prose gate read a literal " as a quote
% character, while \" is LaTeX's diaeresis accent and fails to compile
% inside a node. Every identifier below is written bare, and Speaker:3 is
% typed as it appears in the chapter's prose, not as a quoted string.
%
% Bare tikzpicture (house style): the figure environment, caption and label
% live at the call site in the section file.
\begin{tikzpicture}[
  x=1mm, y=1mm,
  every node/.append style={font=\sffamily\scriptsize},
  posbox/.style={draw, line width=0.4pt, align=center, text width=40mm,
    inner sep=2mm, minimum height=20mm},
  rowlabel/.style={anchor=west, font=\sffamily\scriptsize\itshape},
  crosslink/.style={draw, line width=0.4pt,
    -{Straight Barb[length=1.4mm, width=1.6mm]}},
  droplink/.style={draw, line width=0.4pt, dashed},
  span/.style={draw, line width=0.4pt},
  spanlabel/.style={anchor=north, align=center, inner sep=2pt},
  note/.style={anchor=west, align=left, inner sep=2pt},
]

% ------------------------------------------------------------------- row 1
\node[rowlabel] at (0,16)
  {1. what the router asked: session 4, then session 1, in that order};
\node[posbox] (r1c0) at (45,0)  {position 0\\Speaker:3};
\node[posbox] (r1c1) at (105,0) {position 1\\Speaker:1};

% ------------------------------------------------------------------- row 2
\node[rowlabel] at (0,-26) {2. what the store returned, in its own order};
% The statement belongs to the row rather than to either box, and the row
% label cannot carry it: at this font a label reaching past x=45 runs into
% the left crossing line. Parked right of x=110 instead, which is clear of
% both crossing lines and inside the 155mm measure.
%
% Named rather than written as SQL on purpose. The book's StatementLog
% prints parameter placeholders, so every statement chapters 3, 9 and 10
% print reads IN (@ids1, @ids2); a figure showing IN (3, 1) would be SQL
% no capture in this book produces. What the figure needs is which keys
% went out and in what order, and that is what this says.
\node[note, anchor=east] at (150,-26)
  {one statement, asking for keys 3 and 1};
\node[posbox] (r2c0) at (45,-42)  {row 0\\Ada Fischer, id 1};
\node[posbox] (r2c1) at (105,-42) {row 1\\Chidi Okafor, id 3};

% ------------------------------------------------------------------- row 3
\node[rowlabel] at (0,-64) {3. what the subgraph answered};
\node[posbox] (r3c0) at (45,-80)  {position 0\\Ada Fischer};
\node[posbox] (r3c1) at (105,-80) {position 1\\Chidi Okafor};

% ------------------------------------------------------------------- row 4
\node[rowlabel] at (0,-102) {4. what the client was told};
\node[posbox] (r4c0) at (45,-118)
  {position 0\\The Bank That Split\\Its Graph, by\\Ada Fischer};
\node[posbox] (r4c1) at (105,-118)
  {position 1\\Schemas That Outlive\\Their Authors, by\\Chidi Okafor};

% ------------------------------------------------- the crossing: row 1 to 2
% The whole event. Speaker:3 asked at position 0 is truly answered by row 1
% (Chidi Okafor), and Speaker:1 asked at position 1 is truly answered by
% row 0 (Ada Fischer): the ascending-id order of the store is the reverse
% of the client order the router asked in, so these two lines cross.
\draw[crosslink] (r1c0.south) -- (r2c1.north);
\draw[crosslink] (r1c1.south) -- (r2c0.north);

% -------------------------------------------- straight drops: row 2 to 3
% No crossing from here on: each slot is carried forward by its own
% position, never by the identity actually sitting in it.
\draw[droplink] (r2c0.south) -- (r3c0.north);
\draw[droplink] (r2c1.south) -- (r3c1.north);

% -------------------------------------------- straight drops: row 3 to 4
\draw[droplink] (r3c0.south) -- (r4c0.north);
\draw[droplink] (r3c1.south) -- (r4c1.north);

% ------------------------------------------------------------------ brace
\draw[span] (24,-132) -- (24,-138) -- (126,-138) -- (126,-132);
\node[spanlabel, text width=115mm] at (75,-140)
  {nothing between the store and the client compares a returned entity
   with the representation that asked for it};

% ----------------------------------------------------------------- legend
\draw[crosslink] (4,-152) -- (14,-152);
\node[note, text width=100mm] at (18,-152)
  {the true match: crosses because the store answered in its own order};

\draw[droplink] (4,-160) -- (14,-160);
\node[note, text width=100mm] at (18,-160)
  {carried forward by position alone, never checked again};

\end{tikzpicture}

  \caption{Where the identity goes. The client's order becomes the router's
  representation order, and one \code{IN} list later the store has handed the
  rows back in its own order instead, which is the crossing. Everything below
  it is faithful: the loader writes the row order straight into the position
  order, the subgraph preserves that, the router preserves it, and all three
  are preserving a position that stopped meaning anything two steps earlier.
  What the client is told at the bottom is well-formed, positionally intact,
  and about two real people who gave two other talks. Only the crossing is
  drawn solid, because it is the only link anybody asserted.}
  \label{fig:ch11-position-is-the-join}
\end{figure}

\subsection{GreenDonut Says It Too}

\index{GreenDonut!the contract in the method you override}%
The interesting part is that this graph did not have to go looking for the
specification to get it right. It was told, in the file the defect is in, by
the library chapter~\ref{ch:router-n-plus-1} met when it went looking for a
key-preserving projection. \code{DataLoaderBase.FetchAsync} carries a
documentation comment, and a reader who types \code{override} in an editor is
shown it before writing a line:

\begin{minted}[fontsize=\footnotesize]{text}
A batch loading function which has to be implemented for each
individual DataLoader. For every provided key must be a result
returned. Also to be mentioned is, the results must be returned
in the exact same order the keys were provided.
\end{minted}

\index{DataLoader!three statements of one rule}%
The parameter and the return value each carry their own version of the same
sentence, so the rule is stated three times in one comment. It is the
specification's requirement restated one layer down, which makes sense: a
positional span in and a positional span out is the same contract
\code{\_entities} has, in a method that knows nothing about federation.

\index{DataLoader!the contract has two halves}%
Read it as two clauses rather than one. \emph{For every provided key must be a
result returned} is the first, \emph{in the exact same order the keys were
provided} is the second, and what it costs to break each of them is not the
same thing twice.

\subsection{This Is Not the Guard on the Key}

\index{reference resolver!two different wrong entities}%
Chapter~\ref{ch:first-subgraph} put a type-name check in front of the decoded
key and showed what the entity route answers without one: a \code{Speaker}
node id under a \code{Session} type name comes back as the session carrying
that number, correctly typed, at 200, with no error. That is also right data
and the wrong entity, and it is a different mechanism entirely.

\index{reference resolver!the guard is still running}%
There, the resolver was handed a question it should have refused and answered
it anyway. The fix was to read the type name inside the key. Here the resolver
is handed a question it is entitled to answer, asks for exactly the right
speaker, and is given the wrong one on the way back. The guard
chapter~\ref{ch:first-subgraph} wrote is still in the file, still running, and
still correct, and it is no help at all: it inspects the key going in, and
nothing in this book inspects what comes out.
\index{entities@\code{\_entities}!the order clause, twice|)}%
